% =============================================================================
% COMO LER ESTE LIVRO — PROTOCOLO DO ARQUIVO
% =============================================================================
\cleardoublepage
\thispagestyle{empty}
\chapter*{Como Ler Este Livro}
\addcontentsline{toc}{chapter}{Como Ler Este Livro}
\markboth{Como Ler Este Livro}{}

\begingroup

Este arquivo é um hipertexto impresso de \textbf{84 fragmentos} e \textbf{216 rotas de leitura}. Cada fragmento é uma unidade autônoma de memória ou prova. Você pode percorrê-lo de quatro formas distintas:

\vspace{0.8em}

\noindent\textbf{1. ROTA LINEAR (Cronológica)}\\
\noindent Siga a sequência das páginas, da infância no Sertão (1914) ao Epílogo. É o caminho recomendado para a primeira leitura: acompanha a evolução de Joaquim e a propagação da entidade em ordem temporal.

\vspace{0.5em}

\noindent\textbf{2. ROTA INVESTIGATIVA (Forense)}\\
\noindent Inicie pela \textbf{Parte 4 (Investigação Lúcia, 2026)}. Descubra as anomalias do diário junto com a psiquiatra forense e retorne às Partes 1, 2 e 3 para remontar as provas como uma autópsia.

\vspace{0.5em}

\noindent\textbf{3. ROTA HIPERTEXTUAL (Navegação Livre)}\\
\noindent Ao final de cada fragmento, siga as setas de indicação (\textbf{↪ Rotas}) para saltar diretamente entre as páginas sugeridas. Cada decisão de salto reflete a própria fragmentação da mente do Paciente 1.260. Nenhum caminho é um beco: de qualquer fragmento existe percurso até o Epílogo.

\vspace{0.5em}

\noindent\textbf{4. ROTAS CURADAS (Imersão Dirigida)}\\
\noindent Sequências pré-selecionadas para experiências específicas de contágio:\\[0.3em]
\noindent\textcolor{sangue}{⟡ Rota do Terror:} \rota{CONT-03} $\rightarrow$ \rota{CONT-04} $\rightarrow$ \rota{MEM-12} $\rightarrow$ \rota{CONT-05} --- a propagação da entidade no leitor.\\[0.2em]
\noindent\textcolor{sangue}{⟡ Rota da Compaixão:} \rota{MEM-01} $\rightarrow$ \rota{MEM-05} $\rightarrow$ \rota{MEM-10} $\rightarrow$ \rota{MEM-15} --- a travessia de Joaquim.\\[0.2em]
\noindent\textcolor{sangue}{⟡ Rota Forense:} \rota{DOC-01} $\rightarrow$ \rota{LUC-01} $\rightarrow$ \rota{DOC-16} $\rightarrow$ \rota{CONT-05} --- a investigação do ciclo.

\vspace{0.8em}

\noindent\rule{\textwidth}{0.5pt}
\vspace{0.6em}

\noindent\textbf{CÓDIGOS DOS FRAGMENTOS}

\vspace{0.3em}
\noindent\textcolor{mem}{\textbf{MEM}} --- Memórias do Paciente 1.260 (Sertão e Colônia)\\[0.25em]
\textcolor{doc}{\textbf{DOC}} --- Documentos oficiais, laudos periciais e cartas\\[0.25em]
\textcolor{luc}{\textbf{LUC}} --- Prontuário de investigação da Dra. Lúcia Mendes (2026)\\[0.25em]
\textcolor{cont}{\textbf{CONT}} --- Protocolos de contaminação e resposta psicossomática do leitor

\vspace{1.2em}
\begin{center}\itshape
Nenhum caminho é neutro.\\[0.2em]
Toda leitura é um contágio.\\[0.6em]
\upshape Ao final, você responderá pelo código \textbf{Paciente 1.263}.
\end{center}

\endgroup
\clearpage
